% Generated draft: requires mathematical and source review.
\documentclass[11pt]{article}
\usepackage[margin=25mm]{geometry}
\usepackage{fontspec}
\setmainfont{Latin Modern Roman}
\usepackage{amsmath,amssymb,mathtools,unicode-math}
\setmathfont{Latin Modern Math}
\usepackage{enumitem,needspace,xurl,hyperref}
\setlength{\parindent}{0pt}
\setlength{\parskip}{5pt}
\setlength{\emergencystretch}{3em}
\newcommand{\N}{\mathbb{N}}
\newcommand{\Z}{\mathbb{Z}}
\newcommand{\Q}{\mathbb{Q}}
\newcommand{\R}{\mathbb{R}}
\newcommand{\C}{\mathbb{C}}
\newcommand{\F}{\mathbb{F}}
\newcommand{\A}{\mathbb{A}}
\newcommand{\PP}{\mathbb{P}}
\newcommand{\E}{\mathbb{E}}
\newcommand{\eps}{\varepsilon}
\newcommand{\dd}{\,\mathrm{d}}
\begin{document}
\section*{UnsolvedMath ID: 30004165}
\textbf{Source problem code:} catalogue entry
\subsection{Definitions and mathematical statement}
A locally compact abelian group is an abelian topological group whose topology is Hausdorff and locally compact. An ordinary Haar measure on such a group \(G\) is a nonzero positive regular Borel measure \(\mu\), finite on compact sets, that satisfies
\[
\mu(E+g)=\mu(E)
\qquad\text{for every Borel set }E\subseteq G\text{ and every }g\in G.
\]
This definition supplies terminology only; it does not specify the requested equivariant object.

\emph{Source formulation request.} For locally compact abelian groups \(G\) with a right action by a finite-dimensional semisimple algebra \(A\), what is an appropriate notion of an \emph{equivariant Haar measure}?

\emph{Unresolved qualifications.} The supplied record does not specify the base field of \(A\), its topology, whether the right action means a unital module structure, or any continuity requirement on that action. It also does not specify the values of the proposed measure or the transformation law meant by equivariance. Consequently, the record supports a request for a definition, but does not determine a precise existence, uniqueness, or classification assertion. The ordinary Haar measure definition above does not resolve these missing choices.

The supplied dated literature assessment further states that the extracted question was only ``What does this even mean?'' and that the surrounding prose did not determine a mathematical claim. This qualification remains unresolved despite the separate supplied metadata marking the expanded statement as \texttt{corrected\_verified}.
\subsection{Short English statement}
The record asks how to define an equivariant version of Haar measure for a locally compact abelian group carrying a right algebra action. The algebra, action, and intended equivariance law need further specification. This is a definition request requiring human mathematical review, with no determinate theorem supplied.
\subsection{Sources}
\begin{itemize}
\item Supplied local registry record, \texttt{data/unsolvedmath/problems.json}, permanent UnsolvedMath ID \(30004165\), source problem code \texttt{catalogue entry}, title: \emph{Equivariant Haar Measures for Algebra Actions}.
\item Supplied report citation: \emph{Algebraic K-theory (2020)}, Oberwolfach Reports, pp.\ 1737--1790, DOI \texttt{10.4171/owr/2019/29}. Supplied source URL and literature reference labeled ``OWR primary report'': \url{https://doi.org/10.4171/owr/2019/29}. No individual proposer is identified; the registry gives the proposed year as 2019.
\item Literature assessment dated exactly \texttt{2026-08-22}: the supplied curation labels the record open, but describes the extracted question as ``What does this even mean?'' and the surrounding commentary as explanatory prose without a determinate mathematical claim. This is supplied dated curation, not an independently verified literature assessment. The supplied background requests independent verification of time-sensitive status and cited claims.
\item The supplied statement-verification metadata reports manual semantic, typographic, and LaTeX review and labels the expanded statement \texttt{corrected\_verified}. This is a supplied review claim; its tension with the dated literature assessment requires review of the primary report.
\end{itemize}
\end{document}
